\section{Introduction}

Model validation in chemometrics depends not only on the fitted algorithm, but also on whether the test data represent the population and prediction problem of interest. A numerical score can appear precise while answering an unintended validation question if the partitioning rule changes chemical coverage, response distributions, or the relation between training and test observations. Recent guidance on data-driven model validation likewise emphasizes independent test construction, transparent validation plans, and explicit reporting of their limitations \citep{camacho2026validation}.

Molecular property prediction provides a useful setting in which to study this problem. Public resources such as MoleculeNet and Therapeutics Data Commons have enabled standardized comparisons across chemical datasets, representations, algorithms, and metrics \citep{wu2018moleculenet,huang2021tdc}. However, the resulting benchmark score remains conditional on the split protocol used to define generalization. Random splitting often places chemically similar molecules in both training and test sets. Scaffold splitting attempts to evaluate transfer to held-out molecular frameworks by grouping molecules according to Bemis--Murcko scaffolds \citep{bemis1996properties}, but scaffold disjointness alone does not ensure that the resulting train and test samples are otherwise comparable.

A scaffold partition can alter label prevalence or continuous-property distributions, concentrate the test set in a small number of scaffold groups, and still retain chemically similar molecules across nominally different scaffolds. Consequently, a performance difference between random and scaffold evaluation can combine several effects: reduced local chemical similarity, response-distribution shift, scaffold concentration, and model-specific sensitivity to the selected test region. These concerns connect molecular benchmark design with established QSAR and QSPR principles of external validation and applicability-domain assessment \citep{tropsha2003earnest,gramatica2007principles,netzeva2005applicability,sahigara2012applicability}.

This study treats sample partitioning as an object of chemometric diagnosis. We compare random, ordinary scaffold, and target-balanced scaffold protocols across six public molecular property datasets. The target-balanced protocol preserves scaffold-level separation while searching for a test set whose size and target mean are close to prespecified values. It is used as a diagnostic counterfactual rather than a proposed universal standard: by reducing one measurable source of shift, it tests whether the apparent scaffold penalty persists, shrinks, or changes direction.

The work makes four contributions. First, the target-balanced search is stated as an explicit objective with reproducible constraints and stochastic search settings. Second, its achieved balance is compared with a 5,000-draw random-scaffold null distribution. Third, model effects are evaluated over 20 fixed seeds using paired effect estimates, bootstrap confidence intervals, Wilcoxon signed-rank tests, and Holm correction. Fourth, predictive scores are interpreted jointly with target shift, scaffold composition, test-to-train Tanimoto similarity, sensitivity cases, and model-ranking stability. The purpose is not to introduce a new molecular predictor or to declare one partition universally superior. The purpose is to show which conclusions remain stable after the split itself is measured.